\documentclass[10pt,a4paper]{amsart}
\usepackage[margin=19mm]{geometry}
\usepackage{amsmath,amssymb,mathtools}
\usepackage[scr=rsfs,cal=euler]{mathalfa}
\usepackage{xfrac}
\usepackage[unicode,hidelinks,bookmarks=false]{hyperref}
\newcommand{\ie}{i.e.\ }
\newcommand{\cf}{cf.\ }
\newcommand{\resp}{respectively\ }
\newcommand{\Subsection}{\subsection}
\providecommand{\nobreakdash}{\nobreak\hbox{-}}
\setlength{\emergencystretch}{2em}
\multlinegap0pt
\usepackage{braket}
\usepackage{amssymb}
\usepackage{float}
\usepackage{xcolor}
\usepackage{algorithm}\usepackage{algpseudocode}
\usepackage{tikz}
\usepackage{csquotes}
\newtheorem{theorem}{Theorem}
\title{\textbf{M-Polynomial of Product Graphs}}
\author{El-Mehdi Mehiri$^{a, }$ \and Sandi Klav{\v{z}}ar$^{b, c, d, }$}
\date{Source: arXiv:2603.10596v1 (CC BY 4.0)}
\begin{document}
\maketitle

\noindent\emph{Source and licence.} \textbf{M-Polynomial of Product Graphs}. Version arXiv:2603.10596v1 (CC BY 4.0), distributed by arXiv under the Creative Commons Attribution 4.0 International licence, http://creativecommons.org/licenses/by/4.0/, as recorded in the arXiv licence metadata for this exact version and shown on the abstract page https://arxiv.org/abs/2603.10596v1. Full licence notice: \texttt{licenses/arxiv-2603.10596-CC-BY-4.0.txt}.\par\smallskip

\noindent\textit{Editorial scope.} Counted unit: the article's theorem prop:M-disjunction (source0.tex lines 1457--1488). The article's complete original proof of it is delivered with it (source0.tex lines 1490--1552, one \texttt{proof} environment of 63 lines). No mathematics is written, repaired or re-derived: the statement and the proof are the article's own lines, in the article's own notation, and no retained line is substituted or re-flowed.  No further source-local context is needed: every cross-reference of every retained line resolves inside the two delivered spans.  Nothing was omitted from any retained span, so the package carries no omission record.  No retained line contains a citation, so the wrapper carries no bibliography and no bibliography entry.  No retained line contains a figure environment, an image-inclusion command or drawing code, so no image is delivered and the package declares no asset dependency.  Editorial formatting only, in addition: an AMS \texttt{amsart} preamble that restates the article's own declarations and macros used by the retained lines, namely the environments \texttt{theorem} on separate counters, together with the standard packages those lines need.  The retained environments print in this document as Theorem 1 in order of appearance, whereas the article prints the same items with the numbering of its own full document.  The complete native source of the article as distributed in the arXiv e-print for this exact version is bundled unchanged as \texttt{sources/196148/source0.tex}.
\begin{theorem}
\label{prop:M-disjunction}
If $G$ and $H$ are graphs, then 
\begin{equation}\label{eq:MPolyDisjunction}
    \begin{aligned}
M(G\vee H;x,y)
={}&
\sum_{i\le i'}\ \sum_{j,j'}
m_{i,i'}(G) n_j(H) n_{j'}(H) 
x^{\min\{\delta_\vee(i,j), \delta_\vee(i',j')\}}
y^{\max\{\delta_\vee(i,j), \delta_\vee(i',j')\}}
\\
&+
\sum_{j\le j'}\ \sum_{i,i'}
m_{j,j'}(H) n_i(G) n_{i'}(G) 
x^{\min\{\delta_\vee(i,j), \delta_\vee(i',j')\}}
y^{\max\{\delta_\vee(i,j), \delta_\vee(i',j')\}}
\\
&-
\sum_{i\le i'}\ \sum_{j\le j'}
m_{i,i'}(G) m_{j,j'}(H) 
\Bigl(
x^{\min\{\delta_\vee(i,j), \delta_\vee(i',j')\}}
y^{\max\{\delta_\vee(i,j), \delta_\vee(i',j')\}}
\\
&\hspace{50mm}+
x^{\min\{\delta_\vee(i,j'), \delta_\vee(i',j)\}}
y^{\max\{\delta_\vee(i,j'), \delta_\vee(i',j)\}}
\Bigr).
\end{aligned}
\end{equation}
\end{theorem}

\begin{proof}
By the definition of $G\vee H$ we have $E(G\vee H)=E_A\cup E_B$, 
where
\begin{align*}
E_A & = \bigl\{\{(g,h),(g',h')\}:\ gg'\in E(G)\bigr\}\,, \\
E_B & = \bigl\{\{(g,h),(g',h')\}:\ hh'\in E(H)\bigr\}\,.
\end{align*}
Moreover,
\[
E_A\cap E_B
=\bigl\{\{(g,h),(g',h')\}:\ gg'\in E(G), hh'\in E(H)\bigr\}\,.
\]
Fix now an edge $gg'\in E(G)$ whose endpoint degrees in $G$ are $(i,i')$ (with $i\le i'$).
In $E_A$, the second coordinates are arbitrary, so for any choice of
$h,h'\in V(H)$ with $\deg_H(h)=j$ and $\deg_H(h')=j'$, the unordered pair
$\{(g,h),(g',h')\}$ is an edge of $E_A$.
The number of such choices is $n_j(H) n_{j'}(H)$.
For each such edge, the endpoint degrees in $G\vee H$ are
$\delta_\vee(i,j)$ and $\delta_\vee(i',j')$, hence its M-monomial is
\[
x^{\min\{\delta_\vee(i,j),\delta_\vee(i',j')\}}
y^{\max\{\delta_\vee(i,j),\delta_\vee(i',j')\}}.
\]
Summing over all edges of $G$ of type $(i,i')$, of which there are $m_{i,i'}(G)$,
and over all $(j,j')$, gives exactly the first double sum in the statement, which is
$M_A=\sum_{e\in E_A}x^{\min}y^{\max}$. By the symmetry, fixing an edge $hh'\in E(H)$ yields the second sum. 

Consider next the contribution of the intersection $E_A\cap E_B$. 
Fix an edge $gg'\in E(G)$ of type $(i,i')$ and an edge $hh'\in E(H)$ of type $(j,j')$.
There are \emph{exactly two} distinct edges in $E_A\cap E_B$ generated by these two edges:
\[
\{(g,h),(g',h')\}
\quad\text{and}\quad
\{(g,h'),(g',h)\}.
\]
(They are distinct because $h\neq h'$ and $g\neq g'$.)
Their endpoint degree pairs in $G\vee H$ are
\[
\bigl(\delta_\vee(i,j),\delta_\vee(i',j')\bigr)
\quad\text{and}\quad
\bigl(\delta_\vee(i,j'),\delta_\vee(i',j)\bigr),
\]
hence their $M$-monomials are the two terms inside the parentheses in the third sum.
There are $m_{i,i'}(G)m_{j,j'}(H)$ ways to pick the underlying edges $gg'$ and $hh'$,
so summing over all $(i,i')$ and $(j,j')$ yields
\[
M_I=\sum_{e\in E_A\cap E_B}x^{\min}y^{\max}\,,
\]
which is equal to the third (subtracted) sum in the statement.

Since $E(G\vee H)=E_A\cup E_B$ and $E_A,E_B$ are finite,
\[
\sum_{e\in E(G\vee H)}x^{\min}y^{\max}
=
\sum_{e\in E_A}x^{\min}y^{\max}
+
\sum_{e\in E_B}x^{\min}y^{\max}
-
\sum_{e\in E_A\cap E_B}x^{\min}y^{\max},
\]
that is, $M(G\vee H;x,y)=M_A+M_B-M_I$, which is precisely the formula claimed.
\qedhere
\end{proof}

\end{document}
